%% ============================================================
%% Cross-Database Generalisation of ML-Driven Replication
%% Mode Selection — Journal of Grid Computing (Springer)
%% Template: Springer Nature sn-jnl.cls (sn-mathphys, Numbered)
%% Compiler : pdflatex
%% Status   : DRAFT — results-grounded, not yet reviewed by
%%            advisor. Remaining coordination notes marked [TODO].
%% ============================================================

\documentclass[pdflatex,sn-mathphys-num]{sn-jnl}

% ---- Standard packages (per Journal of Grid Computing guidelines) ----
\usepackage{graphicx}
\usepackage{multirow}
\usepackage{amsmath,amssymb,amsfonts}
\usepackage{amsthm}
\usepackage{manyfoot}
\usepackage{booktabs}
\usepackage{xcolor}
\usepackage{textcomp}
\usepackage{url}
\usepackage{xurl}

% ---- Theorem-style environments (unused here, kept for consistency) ----
\theoremstyle{thmstyleone}
\newtheorem{theorem}{Theorem}
\newtheorem{proposition}[theorem]{Proposition}
\theoremstyle{thmstylethree}
\newtheorem{definition}{Definition}

\raggedbottom

%% ============================================================
\begin{document}

\title[Cross-Database Generalisation of Replication Mode Selection]{%
Cross-Database Generalisation of Machine Learning-Driven
Replication Mode Selection: A Comparative Study Across PostgreSQL,
MySQL, and MongoDB}

\author*[1]{\fnm{Imane} \sur{Jerrari}}\email{imane.jerrari-etu@etu.univh2c.ma}

\author[1]{\fnm{Ismail} \sur{Assayad}}\email{ismail.assayad@univh2c.ma}

\affil*[1]{\orgdiv{Laboratory of Information Systems (LIS)},
\orgname{Faculty of Sciences Ain Chock \& ENSEM, Hassan II
University of Casablanca},
\orgaddress{\city{Casablanca}, \postcode{20000}, \country{Morocco}}}

\abstract{%
% [TODO: numbers below are drawn directly from the experiments run on
%  dataset_mongodb_200.csv, dataset_mysql_200.csv, dataset_postgresql_644.csv.
%  Revise wording once advisor confirms final framing. Word count
%  trimmed to ~230 words to satisfy the journal's 150--250-word limit
%  -- recheck count if further edited.]
Our prior work introduced IAHA, a hybrid machine-learning framework
that selects among four replication modes (PERFORMANCE, CONSISTENCY,
AVAILABILITY, ECONOMIC) for containerised PostgreSQL, leaving
generalisation to other engines as future work. This paper addresses
that gap: does a per-engine RandomForest classifier trained on the
same ten operational features generalise across PostgreSQL, MySQL,
and MongoDB, or does performance depend on engine-specific factors
such as sample size and class balance? We train independent
classifiers on three datasets (PostgreSQL, $n=644$; MySQL, $n=200$;
MongoDB, $n=200$) sharing an identical, engine-agnostic feature
schema. Using repeated stratified 5-fold cross-validation, all three
classifiers significantly outperform a static CPU-threshold heuristic
(McNemar test, $p<10^{-7}$; $+18.0$ to $+21.0$ percentage points).
Mean accuracy is $97.55\%$ for MongoDB, $95.95\%$ for MySQL, and
$98.51\%$ for PostgreSQL. Because PostgreSQL's dataset is $3.2\times$
larger, we test whether its advantage reflects engine identity or
sample size: a size-matched control (PostgreSQL subsampled to
$n=200$) shows accuracy drops by $4.24$ points and falls
\emph{below} both other engines, confirming that dataset size ---
not engine identity --- dominates the observed differences.
Feature-importance rankings correlate strongly across engines
(Spearman $\rho = 0.81$--$0.92$), showing the same operational signals
matter in the same order regardless of engine, a necessary but not
sufficient precondition for transfer learning. However, all three
classifiers fail identically on AVAILABILITY (recall $=0\%$ in every
engine), which we show is a structural feature-space overlap with
CONSISTENCY rather than a correctable sample-size artefact (SMOTE on
PostgreSQL improves recall only to $20\%$). We discuss
the implications for cross-database transfer learning, identified as
future work in our prior study.}

\keywords{High Availability, Containerised Databases, Cross-Database
Generalisation, Machine Learning, Replication Mode Selection}

\maketitle

%% ============================================================
\section{Introduction}\label{sec:intro}
%% ============================================================

\subsection{Context and Motivation}\label{subsec:context}

% [TODO: adjust citation key \cite{iaha2026} to the final reference
%  once the companion paper is accepted / has a DOI]
Adaptive replication-mode selection has recently been shown to
outperform static heuristics on containerised PostgreSQL clusters
\cite{iaha2026}, using a RandomForest classifier trained on operational
metrics to choose among four modes at runtime. That study, however,
was scoped explicitly to a single database engine, and identified
transfer learning to MySQL and MongoDB as a direct item of future
work (FW4 in \cite{iaha2026}).

This leaves an open empirical question that matters beyond PostgreSQL:
is the relationship between operational metrics (CPU utilisation,
transaction throughput, replication lag, etc.) and the appropriate
replication mode a property of PostgreSQL specifically, or a more
general property of replicated database systems under load? If the
latter, a single engine-agnostic feature schema could support adaptive
high availability (HA) across heterogeneous database fleets --- a
common scenario in practice, where organisations operate relational
and NoSQL engines side by side.

\subsection{Problem Statement}\label{subsec:problem}

Does a machine-learning classifier for replication mode
selection, trained independently on PostgreSQL, MySQL, and MongoDB
using an identical engine-agnostic feature schema, achieve consistent
classification performance across engines, and do the same classes
remain difficult to predict regardless of the underlying engine?

\subsection{Research Questions}\label{subsec:rqs}

\begin{itemize}
\item[\textbf{RQ1}] Does classifier performance (accuracy, macro-F1)
  generalise consistently across PostgreSQL, MySQL, and MongoDB, or
  does it depend on engine-specific factors such as dataset size and
  class distribution?
\item[\textbf{RQ2}] Is the relative importance of operational features
  for mode selection stable across a relational engine (PostgreSQL),
  a second relational engine (MySQL), and a document-oriented engine
  (MongoDB)?
\item[\textbf{RQ3}] Is the AVAILABILITY-class failure mode identified
  as a limitation in \cite{iaha2026} (L3: extreme class imbalance)
  specific to PostgreSQL, or does it persist structurally across
  engines regardless of mitigation (e.g.\ oversampling)?
\end{itemize}

\subsection{Contributions}\label{subsec:contributions}

\begin{itemize}
\item An empirical cross-database comparison of replication-mode
  classification across three production-grade database engines under
  an identical, engine-agnostic feature schema.
\item A sample-size control experiment that directly disentangles
  engine identity from training-set size, showing that the apparent
  performance advantage of the largest dataset (PostgreSQL) is
  substantially --- though not entirely --- attributable to its size
  rather than to properties of the engine itself.
\item Statistically grounded evidence (BCa bootstrap confidence
  intervals, McNemar tests against a static heuristic baseline) that
  ML-driven mode selection significantly outperforms static thresholds
  on all three engines, not only PostgreSQL.
\item A diagnostic analysis on PostgreSQL (the only engine with enough
  AVAILABILITY samples to test) showing that the AVAILABILITY-class
  failure persists after SMOTE oversampling rather than resolving as a
  simple data-volume artefact, combined with consistent feature-space
  overlap between AVAILABILITY and CONSISTENCY observed descriptively
  across all three engines.
\item Evidence of cross-engine stability in feature importance rankings
  (Spearman $\rho \geq 0.81$), motivating future transfer-learning work.
\end{itemize}

%% ============================================================
\section{Related Work}\label{sec:related}
%% ============================================================

\subsection{Adaptive High Availability for Databases}\label{subsec:ha}

Existing PostgreSQL HA tools follow a common pattern: a single
synchronous or asynchronous replication policy chosen once at
deployment time and left unchanged at runtime. Patroni \cite{patroni}
and repmgr \cite{repmgr} provide consensus-based leader election but
do not adapt replication behaviour to load; Stolon \cite{stolon}
follows the same fixed-policy pattern. CloudNativePG \cite{cloudnativepg},
the CNCF-donated Kubernetes operator used as the substrate
in \cite{iaha2026}, improves failover latency through a native
Kubernetes control loop but likewise applies one static replication
configuration for the cluster's entire lifetime. None of these tools
offers more than two operational modes, and none adapts mode selection
to observed workload conditions --- the gap that IAHA \cite{iaha2026}
addresses for PostgreSQL, and that this paper extends to MySQL and
MongoDB.

\subsection{Machine Learning for Database Management}\label{subsec:ml4db}

ML-driven database optimisation is an active area spanning several
problem classes: cardinality estimation and query optimisation
(SageDB \cite{kraska2019}, Neo \cite{marcus2019}, Bao \cite{marcus2021}),
automated configuration tuning
(OtterTune \cite{vanaken2017}, CDBTune \cite{zhang2019}), and resource
management (Autopilot \cite{autopilot2020}, FIRM \cite{firm2020}). A
recurring finding in this literature, most directly relevant to our
RQ1, is that ML-driven tuning gains do \emph{not} automatically
transfer across systems or even across versions of the same system:
OtterTune \cite{vanaken2017} explicitly maps unseen workloads to
previously characterised ones because tuning effort on one DBMS does
not reduce tuning effort on another, and LlamaTune \cite{kanellis2022}
reports that the throughput improvement of its tuning approach over a
baseline optimiser narrows substantially when ported from PostgreSQL
v9.6 to v13.6 on certain workloads (e.g.\ from $20.8\%$ to $3.6\%$ on
YCSB-B), despite both being the same engine. Multi-DBMS evaluations
such as GPTuner \cite{lao2024} and UDO \cite{wang2021udo}, which
benchmark across PostgreSQL and MySQL, similarly report
engine-dependent variation in tuning effectiveness rather than uniform
gains. This body of work motivates our explicit empirical test (RQ1,
RQ1b) of whether classification performance generalises across engines
or is confounded by other factors --- in our case, training-set
size --- rather than assuming generalisation holds by default. None of
these systems addresses replication-mode selection specifically, nor
provides operator-facing explainability, which remains the gap
IAHA \cite{iaha2026} and this paper target.

\subsection{Explainable AI in Critical Systems}\label{subsec:xai}

Explainable AI (XAI) addresses the interpretability of ML-driven
decisions \cite{gunning2019} through post-hoc methods such as
SHAP \cite{lundberg2017} and LIME \cite{ribeiro2016}, counterfactual
explanations \cite{wachter2017}, and rule extraction \cite{guidotti2018}.
XAI has been applied to high-stakes domains including
healthcare \cite{holzinger2019}, autonomous vehicles \cite{atakishiyev2021},
and finance \cite{bussmann2020}, but --- as noted in \cite{iaha2026}
--- its application to database HA management remains largely
unexplored. This paper does not introduce new XAI components beyond
those already proposed in \cite{iaha2026}; instead, the feature
importance analysis in Section~\ref{subsec:rq2} serves a
complementary diagnostic role, testing whether the same decision
signature used for explainability on PostgreSQL is structurally stable
enough to remain meaningful on MySQL and MongoDB.

\subsection{Cross-Database and Transfer Learning Approaches}\label{subsec:crossdb}

Despite extensive work on ML for single-DBMS optimisation
(Section~\ref{subsec:ml4db}), explicit cross-engine generalisation
remains comparatively underexplored. Most multi-DBMS evaluations in
the tuning literature (e.g.\ GPTuner \cite{lao2024},
UDO \cite{wang2021udo}) train and evaluate \emph{separate} models per
engine to demonstrate applicability, rather than testing whether a
model or feature representation \emph{transfers} from one engine to
another --- the same design choice we make in this paper (independent
per-engine classifiers, Section~\ref{sec:methodology}), motivated by
the same practical consideration: heterogeneous internal metrics and
configuration spaces across engines make naive model reuse
non-trivial. Industrial efforts to copy or retrain ML models across
database instances exist as systems-engineering patterns, but
typically address operational mechanics of model transfer (e.g.\
avoiding redundant retraining) rather than evaluating whether the
\emph{learned relationships} generalise across heterogeneous engines.
To our knowledge, no prior work evaluates whether an
operational-metrics-to-replication-mode mapping learned on one
database engine is structurally transferable to a different
engine family (relational to relational, or relational to
document-oriented) --- the specific gap this paper addresses
empirically through RQ1--RQ3.

\subsection{Research Gaps Addressed by This Paper}\label{subsec:gaps}

Building on the four gaps identified in \cite{iaha2026} (binary modes,
static thresholds, no explainability, no temporal optimisation), we
identify a fifth gap specific to cross-engine generalisation:

\textbf{Gap 5 -- Untested Cross-Engine Generalisation:} Existing
ML-driven database management systems are evaluated on a single engine
(or, where multiple engines are used, evaluated independently without
testing whether learned relationships transfer between them), leaving
open whether reported performance gains reflect properties of the
engine, the dataset, or the underlying operational-metric relationships
themselves. This paper directly addresses Gap 5 through the
size-matched control analysis in Section~\ref{subsec:rq1b}.

%% ============================================================
\section{Methodology}\label{sec:methodology}
%% ============================================================

\subsection{Datasets}\label{subsec:datasets}

Three datasets were used, sharing an identical ten-feature schema
(Table~\ref{tab:features}) and four-class label space (PERFORMANCE,
CONSISTENCY, AVAILABILITY, ECONOMIC):

\begin{itemize}
\item \textbf{PostgreSQL}: $n = 644$ expert-validated samples
  (the same dataset used to train IAHA's Phase 2 classifier
  in \cite{iaha2026}).
\item \textbf{MySQL}: $n = 200$ samples, collected under the same
  schema for this study.
\item \textbf{MongoDB}: $n = 200$ samples, collected under the same
  schema for this study.
\end{itemize}

\begin{table}[h]
\caption{Shared Feature Schema Across Engines}\label{tab:features}
\begin{tabular}{@{}ll@{}}
\toprule
Feature & Description \\
\midrule
\texttt{replication\_lag\_mb} & Replication lag (MB) \\
\texttt{transactions\_per\_sec} & Transaction throughput \\
\texttt{cpu\_percent} & CPU utilisation (\%) \\
\texttt{latency\_ms} & Query/operation latency \\
\texttt{trans\_rolled\_back} & Rolled-back transaction rate \\
\texttt{deadlocks} & Deadlock count \\
\texttt{hour\_of\_day} & Time-of-day indicator \\
\texttt{memory\_usage\_pct} & Memory utilisation (\%) \\
\texttt{connections\_active} & Active connection count \\
\texttt{wal\_buffers\_full} & Write-ahead/journal buffer saturation \\
\botrule
\end{tabular}
\end{table}

\textbf{Engine-agnostic feature design.} Features were deliberately
chosen as generic operational metrics rather than PostgreSQL-specific
internals, to enable this cross-engine comparison without
re-engineering the feature set per database. We note that
\texttt{wal\_buffers\_full} is measured as a generic write-ahead/journal
buffer saturation proxy; for MongoDB this corresponds functionally to
oplog buffer pressure rather than a literal PostgreSQL WAL buffer, and
we report it under this generic name rather than renaming it
per-engine, to preserve direct cross-engine comparability of the same
underlying measurement axis. Similarly, \texttt{deadlocks} and
\texttt{trans\_rolled\_back} are reported as generic
concurrency-conflict and transaction-failure indicators rather than
engine-native terminology.

\textbf{Synthetic data and class imbalance.} As in \cite{iaha2026}
(Limitation L2), all three datasets are synthetically generated under
expert-informed thresholds rather than collected from production
traces. Class distributions are imbalanced across all three engines,
with CONSISTENCY accounting for $85.0\%$ (MongoDB), $78.5\%$ (MySQL),
and $79.5\%$ (PostgreSQL) of samples, and AVAILABILITY represented by
only $n=1$, $n=3$, and $n=5$ samples respectively
(Table~\ref{tab:class-dist}).

\begin{table}[h]
\caption{Class Distribution by Engine}\label{tab:class-dist}
\begin{tabular}{@{}lrrr@{}}
\toprule
Mode & MongoDB & MySQL & PostgreSQL \\
\midrule
CONSISTENCY  & 170 (85.0\%) & 157 (78.5\%) & 512 (79.5\%) \\
ECONOMIC     & 17 (8.5\%)   & 22 (11.0\%)  & 57 (8.9\%)   \\
PERFORMANCE  & 12 (6.0\%)   & 18 (9.0\%)   & 70 (10.9\%)  \\
AVAILABILITY & 1 (0.5\%)    & 3 (1.5\%)    & 5 (0.8\%)    \\
\midrule
Total ($n$)  & 200 & 200 & 644 \\
\botrule
\end{tabular}
\end{table}

\subsection{Classification Protocol}\label{subsec:protocol}

For each engine, we train an independent RandomForest classifier
(200 trees, balanced class weighting) using repeated stratified 5-fold
cross-validation (10 repeats, 50 fold-level accuracy estimates per
engine), matching the cross-validation protocol used
in \cite{iaha2026}.

\subsection{Statistical Testing}\label{subsec:stats}

\textbf{BCa bootstrap confidence intervals.} Following the
methodology of \cite{iaha2026}, we report bias-corrected and
accelerated (BCa) bootstrap 95\% confidence intervals (10{,}000
resamples) on mean cross-validated accuracy per engine.

\textbf{McNemar test.} To test whether the RandomForest classifier
significantly outperforms a static heuristic, we apply McNemar's test
\emph{within each engine}, comparing paired correct/incorrect
predictions on the same held-out folds. The static heuristic mirrors
the threshold logic of IAHA's P1 layer: \texttt{cpu\_percent}\,$<30 \Rightarrow$
ECONOMIC; $30\leq$\texttt{cpu\_percent}\,$<50 \Rightarrow$ PERFORMANCE;
otherwise CONSISTENCY (the heuristic cannot express AVAILABILITY).
We do \emph{not} apply McNemar across engines (e.g.\ MongoDB vs.\
MySQL), since the two samples are drawn from different observation
spaces and are not paired --- this would violate the test's
matched-pairs assumption.

\textbf{Mann--Whitney U test.} To compare the distributions of
cross-validated accuracy between pairs of engines, we use the
Mann--Whitney U test on the 50 fold-level accuracy estimates per
engine. We note as a limitation that these 50 estimates per engine are
not fully independent (they derive from repeated resampling of the
same underlying dataset), so the resulting $p$-values likely overstate
the true level of statistical confidence; we report them as indicative
rather than definitive evidence of inter-engine difference.

\subsection{Sample-Size Control Protocol}\label{subsec:sizecontrol-protocol}

Because PostgreSQL ($n=644$) is $3.2\times$ larger than MySQL or
MongoDB ($n=200$ each), any raw accuracy comparison between engines is
confounded with a sample-size effect. To isolate the contribution of
engine identity from dataset size, we additionally evaluate a
\emph{size-matched} PostgreSQL condition: 20 independent stratified
subsamples of $n=200$ are drawn from the full PostgreSQL dataset
(class-proportional: 158 CONSISTENCY, 22 PERFORMANCE, 18 ECONOMIC, 2
AVAILABILITY), each evaluated under the same repeated stratified
cross-validation protocol as the full dataset, yielding 40 fold-level
accuracy estimates. This size-matched PostgreSQL condition is then
compared against (i) PostgreSQL at full size, and (ii) MySQL and
MongoDB at their native $n=200$, using the same Mann--Whitney
procedure described above.

%% ============================================================
\section{Results}\label{sec:results}
%% ============================================================

\subsection{RQ1: Cross-Engine Classification Performance}\label{subsec:rq1}

Table~\ref{tab:main-results} summarises classification performance
across the three engines.

\begin{table}[h]
\caption{Classification Performance by Engine (Repeated $10\times5$-Fold CV)}\label{tab:main-results}
\begin{tabular}{@{}lrrr@{}}
\toprule
Metric & MongoDB & MySQL & PostgreSQL \\
\midrule
Mean accuracy & 97.55\% & 95.95\% & 98.51\% \\
BCa 95\% CI & [97.0, 98.1] & [95.2, 96.6] & [98.3, 98.7] \\
Majority baseline & 85.00\% & 78.50\% & 79.50\% \\
Heuristic acc. & 79.00\% & 75.00\% & 79.19\% \\
Macro-F1 (3 cl.)\footnotemark[1] & 94.95\% & 95.41\% & 98.31\% \\
\botrule
\end{tabular}
\footnotetext[1]{CONSISTENCY, ECONOMIC, PERFORMANCE only; AVAILABILITY
excluded from aggregate macro-F1 due to insufficient sample size for
reliable estimation (see Section~\ref{subsec:rq3}).}
\end{table}

All three classifiers substantially exceed both the majority-class
baseline and the static heuristic. We note that for PostgreSQL the
static heuristic ($79.19\%$) slightly underperforms the majority-class
baseline ($79.50\%$): the heuristic actively misclassifies some
CONSISTENCY samples into ECONOMIC or PERFORMANCE when
\texttt{cpu\_percent} crosses its thresholds, whereas a baseline that
always predicts CONSISTENCY incurs no such cost; this is consistent
with the heuristic providing no information about AVAILABILITY
and a CPU threshold that only weakly separates CONSISTENCY from
PERFORMANCE (Section~\ref{subsec:rq3}). The RandomForest classifier
does not exhibit this issue on any engine. McNemar's test confirms the
RandomForest classifier significantly outperforms the static
heuristic on every engine (Table~\ref{tab:mcnemar}). Note that the RF
accuracy reported in Table~\ref{tab:mcnemar} is computed on a single
stratified 5-fold split (required for McNemar's paired-prediction
design) rather than the repeated $10\times5$-fold protocol used for
Table~\ref{tab:main-results}; the two are not directly comparable and
the small numerical difference (e.g.\ $98.9\%$ vs.\ $98.51\%$ for
PostgreSQL) reflects this difference in protocol, not an
inconsistency in the underlying model.

\begin{table}[h]
\caption{McNemar Test: RandomForest vs.\ Static Heuristic (Per Engine)}\label{tab:mcnemar}
\begin{tabular}{@{}lrrrr@{}}
\toprule
Engine & RF acc. & Heur.\ acc. & $\Delta$ (pp) & McNemar $p$ \\
\midrule
MongoDB     & 97.0\% & 79.0\% & $+18.0$ & $2.4\times10^{-8}$ \\
MySQL       & 96.0\% & 75.0\% & $+21.0$ & $4.1\times10^{-9}$ \\
PostgreSQL  & 98.9\% & 79.2\% & $+19.7$ & $8.6\times10^{-26}$ \\
\botrule
\end{tabular}
\end{table}

The magnitude of improvement ($+18.0$ to $+21.0$ percentage points) is
consistent with the $+25.7$ percentage-point improvement over static
heuristics reported for PostgreSQL in \cite{iaha2026}, independently
confirming on two additional engines that ML-driven mode selection
provides a substantial, statistically significant advantage over fixed
thresholds.

\textbf{Engine-level differences.} A Mann--Whitney U test on
fold-level accuracy distributions found all three pairwise comparisons
statistically significant (MongoDB vs.\ MySQL: $p=0.0014$; MongoDB
vs.\ PostgreSQL: $p=4.7\times10^{-4}$; MySQL vs.\ PostgreSQL:
$p=2.2\times10^{-12}$), though as discussed in
Section~\ref{subsec:stats} these $p$-values should be read as
indicative rather than definitive given the non-independence of
repeated-CV estimates. PostgreSQL --- the largest dataset
($n=644$) --- achieves the highest accuracy and the tightest
confidence interval, which raises an immediate confound already
flagged in Section~\ref{sec:methodology}: is PostgreSQL's apparent
advantage a property of the engine, or simply a property of having
$3.2\times$ more training data? We test this directly in
Section~\ref{subsec:rq1b}.

\subsection{RQ1b: Disentangling Engine Identity from Sample Size}\label{subsec:rq1b}

To test whether PostgreSQL's higher accuracy in
Table~\ref{tab:main-results} reflects engine-specific structure or
simply its larger sample size, we apply the size-matched control
described in Section~\ref{subsec:sizecontrol-protocol}: PostgreSQL
subsampled to $n=200$ (class-proportional, 20 repeats) is compared
against PostgreSQL at full size and against MySQL/MongoDB at their
native $n=200$.

\begin{table}[h]
\caption{Sample-Size Control: PostgreSQL Full vs.\ Size-Matched}\label{tab:size-control}
\begin{tabular}{@{}lrr@{}}
\toprule
Condition & Mean acc. & BCa 95\% CI \\
\midrule
PostgreSQL, full ($n=644$)         & 98.51\% & [98.31, 98.67] \\
PostgreSQL, matched ($n=200$)      & 94.27\% & [93.47, 94.95] \\
MySQL, native ($n=200$)            & 95.95\% & [95.15, 96.60] \\
MongoDB, native ($n=200$)          & 97.55\% & [97.00, 98.10] \\
\botrule
\end{tabular}
\end{table}

The result is unambiguous: once PostgreSQL is reduced to the same
sample size as MySQL and MongoDB, its accuracy drops by $4.24$
percentage points (from $98.51\%$ to $94.27\%$), and its confidence
interval no longer overlaps with the full-size estimate (Mann--Whitney
$U=53.0$, $p=8.6\times10^{-15}$). \textbf{This confirms that the
performance advantage attributed to PostgreSQL in
Section~\ref{subsec:rq1} is substantially driven by dataset size,
not by engine identity.}

The per-class breakdown (Table~\ref{tab:size-control-classes})
localises this effect precisely for the PERFORMANCE class (the
sample-size-matched subsample also contains only $n=2$ AVAILABILITY
cases, too few to report a recall figure for that class here; see
Section~\ref{subsec:rq3} for the AVAILABILITY analysis): PERFORMANCE recall collapses from
$97\%$ (full PostgreSQL) to $59\%$ (size-matched PostgreSQL) once only
22 PERFORMANCE examples are available for training instead of 70 ---
closely matching the recall range already observed for the
data-scarce PERFORMANCE class in MySQL and MongoDB. ECONOMIC remains
at $100\%$ regardless of sample size, consistent with its empirical
separability discussed in Section~\ref{subsec:economic-note}.

\begin{table}[h]
\caption{PERFORMANCE-Class Recall vs.\ Training Set Size}\label{tab:size-control-classes}
\begin{tabular}{@{}lrr@{}}
\toprule
Condition & $n_{\text{PERFORMANCE}}$ & Recall \\
\midrule
PostgreSQL, full ($n=644$)         & 70 & 97\% \\
PostgreSQL, size-matched ($n=200$) & 22 & 59\% \\
MySQL ($n=200$)                    & 18 & 72\% \\
MongoDB ($n=200$)                  & 12 & 58\% \\
\botrule
\end{tabular}
\end{table}

Notably, the size-matched PostgreSQL condition does not fully converge
with MySQL and MongoDB: it remains significantly different from both
(vs.\ MySQL: $U=559.5$, $p=2.8\times10^{-4}$; vs.\ MongoDB:
$U=260.5$, $p=1.1\times10^{-9}$), and --- counter-intuitively --- its
mean accuracy ($94.27\%$) is now \emph{below} both MySQL ($95.95\%$)
and MongoDB ($97.55\%$) at the same sample size. We interpret this
residual difference cautiously. Having confirmed (Limitation L1) that
all three datasets are produced by the same generator and the same
random seed, a difference in generation procedure is not a plausible
explanation here. The more likely candidates are sampling variance not
fully captured by our 20-repeat subsampling protocol, or a genuine
finite-sample property of how PostgreSQL's particular scenario mix
happens to shape the four-class decision boundary at $n=200$
specifically. We do not have grounds to attribute the residual to a
property of the PostgreSQL \emph{engine} itself, and flag this as an
open question for future work (Section~\ref{sec:conclusion}, FW4
below) rather than a settled finding.

\textbf{Implication for RQ1.} The evidence supports a revised answer
to RQ1: classifier performance does \emph{not} generalise uniformly
across engines at matched sample sizes, but the dominant driver of the
cross-engine differences observed in Table~\ref{tab:main-results} is
training-set size --- particularly its effect on data-scarce classes
such as PERFORMANCE --- rather than an intrinsic advantage of any
single database engine. This has a direct practical implication: when
deploying ML-driven mode selection on a new engine, the data
collection budget per class (especially for non-dominant modes) likely
matters more than the choice of engine itself.

\subsection{RQ2: Cross-Engine Feature Importance Stability}\label{subsec:rq2}

\begin{table}[h]
\caption{Top-4 Feature Importances by Engine (RandomForest)}\label{tab:feature-importance}
\begin{tabular}{@{}lrrr@{}}
\toprule
Feature & MongoDB & MySQL & PostgreSQL \\
\midrule
\texttt{transactions\_per\_sec} & 0.227 & 0.279 & 0.238 \\
\texttt{cpu\_percent}           & 0.204 & 0.194 & 0.227 \\
\texttt{connections\_active}    & 0.150 & 0.149 & 0.171 \\
\texttt{replication\_lag\_mb}   & 0.122 & 0.135 & 0.158 \\
\botrule
\end{tabular}
\end{table}

The same four features dominate the decision process across all three
engines, in the same relative order. Spearman rank correlation between
full feature-importance rankings (all ten features) is $\rho = 0.806$
(MongoDB vs.\ MySQL), $\rho = 0.915$ (MongoDB vs.\ PostgreSQL), and
$\rho = 0.891$ (MySQL vs.\ PostgreSQL). We emphasise what this
correlation does and
does not show: it indicates that the \emph{relative ranking} of which
features matter is stable across engines, which is a necessary
precondition for transfer learning to be worth attempting (there
would be little reason to transfer a model if the engines relied on
different signals entirely). It is not, by itself, evidence that a
classifier's learned decision boundary transfers --- that would
require directly training on one engine and testing on another, which
we identify as future work (FW1, Section~\ref{sec:conclusion}) rather
than claim here. We therefore describe the decision signature as
\emph{structurally consistent} across engines, and treat genuine
transferability as an open empirical question.

\subsection{A Note on ECONOMIC Separability}\label{subsec:economic-note}

Throughout this paper, the ECONOMIC class is perfectly separable
(recall and precision $=100\%$ on all three engines and at all sample
sizes tested), and our static heuristic baseline reproduces this
separation using a single threshold, $\texttt{cpu\_percent} < 30$
(Section~\ref{subsec:stats}). It is important to be precise about
\emph{why} this threshold works, since the underlying data-generation
process (consulted directly rather than inferred) reveals that
\texttt{cpu\_percent} is not itself the generating rule for ECONOMIC.
ECONOMIC samples are generated under a low-load
``off-peak'' scenario defined jointly by night-time hours and low
transaction throughput; CPU utilisation in that scenario happens to be
drawn from a low range ($10$--$30\%$) as a \emph{consequence} of the
same scenario, not as its defining criterion. The
$\texttt{cpu\_percent}<30$ threshold is therefore an effective
\emph{empirical proxy} for this scenario, not a description of the
causal mechanism that produced the label. We flag this distinction
explicitly because it affects how the perfect separability of
ECONOMIC should be read: it demonstrates that a single operational
metric can capture a multivariate, time-of-day-dependent operating
regime well enough to be perfectly diagnostic on this dataset, which
is a more interesting empirical finding than a merely definitional
threshold would be, but it should not be over-interpreted as evidence
that CPU utilisation alone \emph{causes} or \emph{defines} the
ECONOMIC mode in a deployed system.

\subsection{RQ3: The AVAILABILITY Failure Mode}\label{subsec:rq3}

Across all three engines, AVAILABILITY recall is $0\%$
(Table~\ref{tab:confusion-availability}): every AVAILABILITY sample is
misclassified, and in every case the misclassification is to
CONSISTENCY --- never to ECONOMIC or PERFORMANCE.

\begin{table}[h]
\caption{AVAILABILITY Misclassification Pattern by Engine}\label{tab:confusion-availability}
\begin{tabular}{@{}lrrrr@{}}
\toprule
Engine & $n_{\text{avail}}$ & Predicted CONS. & Recall & Feature overlap\footnotemark[2] \\
\midrule
MongoDB     & 1 & 1/1 & 0\% & n/a\footnotemark[3] \\
MySQL       & 3 & 3/3 & 0\% & 33.2\% \\
PostgreSQL  & 5 & 5/5 & 0\% & 38.9\% \\
\botrule
\end{tabular}
\footnotetext[2]{Mean \% of CONSISTENCY samples falling within the
observed AVAILABILITY feature range, averaged across the ten features.}
\footnotetext[3]{Not statistically interpretable with $n=1$.}
\end{table}

To test whether this failure is a correctable sample-size artefact or
a structural overlap, we applied SMOTE oversampling to the PostgreSQL
training folds (the only engine with enough AVAILABILITY samples,
$n=5$, to make oversampling meaningful). SMOTE improved AVAILABILITY
recall from $0\%$ to only $20\%$ (F1 $=0.29$), with precision dropping
to $50\%$. This causal test was not repeated on MySQL or MongoDB, as
$n=1$ and $n=3$ AVAILABILITY samples are too few to support
oversampling. Combined with the descriptive feature-overlap evidence
that holds across all three engines ($33$--$39\%$ for MySQL and
PostgreSQL; not statistically interpretable for MongoDB at $n=1$), we
conclude that the failure is more consistent with a structural
overlap than with a purely volumetric shortage --- the current
ten-feature schema does not contain a feature that cleanly separates
AVAILABILITY from CONSISTENCY on the one engine where this could be
directly tested, and the same feature-space overlap is visible
wherever it can be measured. This mirrors and extends Limitation L3 of
\cite{iaha2026}
(``CONSISTENCY $=89.4\%$ limits generalisation for ECONOMIC and
AVAILABILITY''), showing the same limitation persists after extension
to two additional engines and is not specific to the original
PostgreSQL dataset.

%% ============================================================
\section{Discussion}\label{sec:discussion}
%% ============================================================

% [TODO: expand once full draft reviewed by advisor.]
\textbf{RQ1.} Classification accuracy is high across all three engines
but is not uniform, and our size-matched control
(Section~\ref{subsec:rq1b}) confirms \emph{empirically} --- rather
than merely suggesting --- that this non-uniformity is driven
primarily by training-set size rather than engine identity: subsampling
PostgreSQL to MySQL/MongoDB's native size erases its accuracy
advantage entirely, with the loss concentrated in the data-scarce
PERFORMANCE class. This is an important distinction for practitioners:
the data collection budget per class (especially for non-dominant
modes such as PERFORMANCE and AVAILABILITY) likely matters more than
which engine is being instrumented. At the same time, the size-matched
PostgreSQL condition does not fully converge with MySQL or MongoDB
accuracy (Section~\ref{subsec:rq1b}), leaving open the
possibility of a smaller, engine- or dataset-specific residual effect
that this study is not positioned to fully explain.

\textbf{RQ2.} The strong cross-engine correlation in feature importance
($\rho \geq 0.81$) suggests that \texttt{transactions\_per\_sec},
\texttt{cpu\_percent}, \texttt{connections\_active}, and
\texttt{replication\_lag\_mb} constitute a largely engine-agnostic
``signature'' of replication-mode state. This is a necessary (though
not sufficient) condition for transfer learning across engines: a
classifier pre-trained on PostgreSQL's larger dataset could plausibly
provide a useful initialisation for MySQL or MongoDB classifiers,
particularly for the data-scarce PERFORMANCE and AVAILABILITY classes.
% [TODO: this is the natural bridge to a future Paper 3 on transfer
%  learning — make this connection explicit once paper 3 scope is set]

\textbf{RQ3.} The fact that AVAILABILITY is \emph{never} confused with
ECONOMIC or PERFORMANCE --- only ever with CONSISTENCY --- across all
three engines is itself informative: it suggests the current feature
schema captures a single dominant axis of variation (roughly,
load/CPU intensity) that separates ECONOMIC and PERFORMANCE well, but
does not capture whatever distinguishes an ``available-but-degraded''
state from a normal CONSISTENCY state. Candidate additional features
for future work include active replica/read-replica counts, health-check
failure rates, or failover-readiness indicators --- none of which are
present in the current ten-feature schema.

\subsection{Limitations}\label{subsec:limitations}

\textbf{L1 -- Synthetic data, single shared generator.} As in
\cite{iaha2026}, all three datasets are synthetically generated rather
than collected from live production traces. We confirm, by direct
inspection of the generation script, that all three datasets
(PostgreSQL, MySQL, MongoDB) are produced by a single shared Python
generator (a common \texttt{generate\_samples(n, system)} function),
called three times with the same random seed and the same
scenario-based logic, varying only the requested sample size $n$ and
an engine-specific parameter $\texttt{system}$. Each sample is drawn
from one of five latent operating scenarios (normal, peak, off-peak,
crisis, degraded), each with its own probability and its own
parameter ranges for the ten operational features; the four-class
\texttt{mode} label is then assigned by a deterministic rule cascade
over the sampled feature values (load thresholds, error rates,
deadlock counts, and time-of-day), not by a separate validation step
applied to each generated row. ``Expert-validated'' in this context
and in \cite{iaha2026} refers to the scenario probabilities,
parameter ranges, and labelling-rule cascade themselves having been
specified by domain expertise, not to an independent post-hoc review
of each individual sample. Two of the ten features
(\texttt{wal\_buffers\_full}, \texttt{trans\_rolled\_back}) are drawn
from engine-specific ranges within the shared generator (e.g.\
\texttt{wal\_buffers\_full} $\sim \mathcal{U}(0,50)$ for PostgreSQL
vs.\ $\mathcal{U}(0,20)$ for MongoDB), consistent with the $2$--$2.5\times$
scale divergence we measured empirically between engines on exactly
these two features (Section~\ref{subsec:datasets}); all remaining
features share identical generating ranges across engines. This
confirms that cross-engine differences reported in this paper
(Sections~\ref{subsec:rq1}--\ref{subsec:rq3}) reflect a single,
consistent generative process rather than independently calibrated,
potentially incomparable, per-engine implementations. This also
narrows the candidate explanations for the unexplained residual gap
identified in Section~\ref{subsec:rq1b} (size-matched PostgreSQL
underperforming MySQL and MongoDB at equal $n$): a generator
discrepancy is no longer a plausible cause, since the same code and
seed produce all three datasets, which strengthens --- rather than
undermines --- the case that the residual reflects genuine
finite-sample structure in how the four-class decision boundary is
shaped at PostgreSQL's particular scenario mix, and we revise FW4
(Section~\ref{sec:conclusion}) accordingly.

\textbf{L2 -- Extreme AVAILABILITY imbalance.} With $n=1$, $3$, and
$5$ AVAILABILITY samples respectively, no engine permits a
statistically reliable estimate of AVAILABILITY recall; we report it
descriptively rather than within aggregate metrics.

\textbf{L3 -- Non-independence of repeated-CV estimates.} Mann--Whitney
comparisons between engines use 50 fold-level accuracy estimates per
engine drawn from repeated resampling of the same underlying dataset;
these are not fully independent observations, so reported $p$-values
likely overstate true statistical confidence.

\textbf{L4 -- No live multi-engine cluster validation.} Unlike
the live Kubernetes validation in \cite{iaha2026} (Phase 1, CloudNativePG),
this study validates classification performance only; end-to-end
decision latency and live failover behaviour on MySQL/MongoDB
clusters remain to be measured.
% [TODO: if/when a live test is run, this becomes the natural
%  Phase-1-equivalent for paper 2]

\subsection{Threats to Validity}\label{subsec:threats}

\textbf{Internal.} The static heuristic baseline reuses the threshold
logic of IAHA's P1 layer rather than an independently tuned baseline
per engine, which may understate what a carefully engine-tuned
heuristic could achieve.

\textbf{Construct.} Feature semantics (e.g.\ \texttt{wal\_buffers\_full}
as an oplog-buffer proxy for MongoDB) are functional analogues rather
than native engine metrics; this trades engine-specific fidelity for
cross-engine comparability, an explicit design choice discussed in
Section~\ref{sec:methodology}.

\textbf{External.} Results are obtained on synthetic datasets of
differing size (200/200/644); generalisation to production traces and
to balanced multi-engine datasets of equal size remains to be
validated (cf.\ L1, L2).

%% ============================================================
\section{Conclusion and Future Work}\label{sec:conclusion}
%% ============================================================

% [TODO: finalise once advisor confirms final framing.]
This study directly addresses FW4 from \cite{iaha2026} by evaluating
whether ML-driven replication mode selection generalises from
PostgreSQL to MySQL and MongoDB. We find that classification accuracy
is consistently high and significantly outperforms static heuristics
on all three engines (McNemar $p<10^{-7}$), and that the underlying
feature-importance signature is highly stable across engines (Spearman
$\rho \geq 0.81$) --- evidence that supports, but does not yet
demonstrate, cross-engine transfer learning. Critically, a size-matched
control experiment shows that the cross-engine accuracy differences
observed at native dataset sizes are substantially driven by
training-set size rather than by intrinsic properties of the database
engine: a result with direct implications for how practitioners should
budget data collection when extending ML-driven HA to new engines. At
the same time, the AVAILABILITY class fails identically (zero recall,
always misclassified as CONSISTENCY) on all three engines, and the
evidence available --- a direct causal test on PostgreSQL plus
consistent descriptive feature overlap elsewhere --- points to a
genuine gap in the current feature schema rather than an artefact
specific to any one dataset.

\textbf{FW1.} Explicit transfer-learning evaluation (pre-train on
PostgreSQL, fine-tune on MySQL/MongoDB), directly testing whether the
feature-importance stability found here translates into sample-efficiency
gains for the data-scarce engines.

\textbf{FW2.} Extend the feature schema with availability-specific
signals (replica health-check failure rate, active read-replica count)
to test whether the structural AVAILABILITY/CONSISTENCY overlap
identified in Section~\ref{subsec:rq3} can be resolved.

\textbf{FW3.} Live multi-engine cluster validation (MySQL Group
Replication, MongoDB replica sets) mirroring the Phase 1 CloudNativePG
validation in \cite{iaha2026}, to measure real end-to-end decision
latency and failover behaviour rather than offline classification
accuracy alone.

\textbf{FW4.} Investigate the residual, unexplained performance gap
between size-matched PostgreSQL and MySQL/MongoDB identified in
Section~\ref{subsec:rq1b}. Having confirmed (Limitation L1) that all
three datasets share a single generator and seed, a generator
discrepancy is no longer a candidate explanation; remaining
hypotheses are sampling variance not fully captured by our 20-repeat
protocol, or genuine finite-sample structure in how the four-class
decision boundary behaves under PostgreSQL's particular scenario mix
at $n=200$ specifically.

%% ============================================================
\backmatter

\bmhead{Acknowledgments}
% [TODO: add acknowledgments if applicable, or state "Not applicable".]
Not applicable.

\section*{Declarations}

% [TODO: confirm exact wording with advisor before submission.]
\begin{itemize}
\item \textbf{Funding:} The authors did not receive support from any
  organization for the submitted work.
\item \textbf{Competing Interests:} The authors have no relevant
  financial or non-financial interests to disclose.
\item \textbf{Ethics approval:} Not applicable.
\item \textbf{Consent to participate:} Not applicable.
\item \textbf{Consent for publication:} Not applicable.
\item \textbf{Availability of data and materials:} The datasets used
  in this study are available from the corresponding author upon
  reasonable request.
  % [TODO: consider depositing datasets in a public repository
  %  (e.g. Zenodo) per Journal of Grid Computing's research data
  %  policy, and replace this statement with a DOI-bearing citation.]
\item \textbf{Code availability:} The analysis code used to produce
  the results in this paper is available from the corresponding
  author upon reasonable request.
\item \textbf{Authors' contributions:} All authors contributed to the
  study conception and design. Material preparation, data collection,
  and analysis were performed by I.~Jerrari. The first draft of the
  manuscript was written by I.~Jerrari and all authors commented on
  previous versions of the manuscript. All authors read and approved
  the final manuscript.
  % [TODO: confirm/adjust contribution statement with co-author.]
\item \textbf{AI-assisted copy editing:} Large Language Model (LLM)
  tools were used for AI-assisted copy editing (readability and style
  improvements to author-generated text) and for exploratory data
  analysis under direct author supervision; all analytical decisions,
  interpretations, and conclusions are the authors' own. No generative
  AI tool is listed as an author.
  % [TODO: confirm this statement reflects the actual workflow and
  %  satisfies Journal of Grid Computing's LLM-use disclosure
  %  requirements before submission.]
\end{itemize}

% ---- Bibliography (numbered, per Journal of Grid Computing style) ----
\begin{thebibliography}{20}

\bibitem{iaha2026}
Jerrari, I., Assayad, I.: IAHA: An intelligent adaptive high
availability framework for containerised PostgreSQL databases using
hybrid machine learning. Journal of Intelligent \& Fuzzy Systems,
under review (2026)

\bibitem{patroni}
Patroni Project: Patroni: A template for PostgreSQL HA.
\url{https://github.com/patroni/patroni}. Accessed Apr 2025

\bibitem{repmgr}
repmgr Project: repmgr: PostgreSQL replication manager.
\url{https://repmgr.org/}. Accessed Apr 2025

\bibitem{stolon}
Sorint.lab: Stolon: PostgreSQL cloud-native HA.
\url{https://github.com/sorintlab/stolon}. Accessed Apr 2025

\bibitem{cloudnativepg}
EnterpriseDB/CNCF: CloudNativePG: Cloud-native PostgreSQL operator.
\url{https://cloudnative-pg.io/}. Accessed Apr 2025

\bibitem{kraska2019}
Kraska, T., et al.: SageDB: A learned database system. In: Proc. CIDR
(2019)

\bibitem{marcus2019}
Marcus, R., Negi, P., Mao, H., Zhang, C., Alizadeh, A., Kraska, T.,
Papaemmanouil, O., Tatbul, N.: Neo: A learned query optimizer. Proc.
VLDB Endow. \textbf{12}(11), 1705--1718 (2019)

\bibitem{marcus2021}
Marcus, R., et al.: Bao: Making learned query optimization practical.
In: Proc. ACM SIGMOD, pp.\ 1275--1288 (2021)

\bibitem{vanaken2017}
Van Aken, D., Pavlo, A., Gordon, G.J., Zhang, B.: Automatic database
management system tuning through large-scale machine learning. In:
Proc. ACM SIGMOD, pp.\ 1009--1024 (2017)

\bibitem{zhang2019}
Zhang, J., et al.: An end-to-end automatic cloud database tuning
system using deep reinforcement learning. In: Proc. ACM SIGMOD, pp.\
415--432 (2019)

\bibitem{autopilot2020}
Rzadca, K., Findeisen, P., Swiderski, J., Zych, P., Broniek, P.,
Kusmierek, J., Nowak, P., Strack, B., Witusowski, P., Hand, S.,
Wilkes, J.: Autopilot: Workload autoscaling at Google. In: Proc.\ 15th
Eur.\ Conf.\ Comput.\ Syst.\ (EuroSys '20), Heraklion, Greece, Article
no.\ 16 (2020). \url{https://doi.org/10.1145/3342195.3387524}

\bibitem{firm2020}
Qiu, H., et al.: FIRM: An intelligent fine-grained resource
management framework for SLO-oriented microservices. In: Proc.\ OSDI,
pp.\ 805--825 (2020)

\bibitem{gunning2019}
Gunning, D., et al.: XAI---Explainable artificial intelligence. Sci.
Robot. \textbf{4}(37), eaay7120 (2019)

\bibitem{lundberg2017}
Lundberg, S.M., Lee, S.-I.: A unified approach to interpreting model
predictions. In: Proc.\ NeurIPS, pp.\ 4765--4774 (2017)

\bibitem{ribeiro2016}
Ribeiro, M.T., Singh, S., Guestrin, C.: ``Why should I trust you?'':
Explaining the predictions of any classifier. In: Proc.\ ACM KDD, pp.\
1135--1144 (2016)

\bibitem{wachter2017}
Wachter, S., Mittelstadt, B., Russell, C.: Counterfactual explanations
without opening the black box: Automated decisions and the GDPR.
Harvard J. Law Technol. \textbf{31}(2), 841--887 (2017)

\bibitem{guidotti2018}
Guidotti, R., et al.: A survey of methods for explaining black box
models. ACM Comput. Surv. \textbf{51}(5), 93:1--93:42 (2018)

\bibitem{holzinger2019}
Holzinger, A., Langs, G., Denk, H., Zatloukal, K., M\"uller, H.:
Causability and explainability of artificial intelligence in
medicine. WIREs Data Min. Knowl. Discov. \textbf{9}(4), e1312 (2019).
\url{https://doi.org/10.1002/widm.1312}

\bibitem{atakishiyev2021}
Atakishiyev, S., Salameh, M., Yao, H., Goebel, R.: Explainable
artificial intelligence for autonomous driving. arXiv:2104.08952
(2021)

\bibitem{bussmann2020}
Bussmann, N., et al.: Explainable machine learning in credit risk
management. Comput. Econ. \textbf{57}, 203--216 (2021)

\bibitem{kanellis2022}
Kanellis, K., Ding, C., Kroth, B., M\"uller, A., Curino, C.,
Venkataraman, S.: LlamaTune: Sample-efficient DBMS configuration
tuning. Proc. VLDB Endow. \textbf{15}(11), 2953--2965 (2022)

\bibitem{lao2024}
Lao, J., Wang, Y., Li, Y., Wang, J., Zhang, Y., Cheng, Z., Chen, W.,
Tang, M., Wang, J.: GPTuner: A manual-reading database tuning system
via GPT-guided Bayesian optimization. Proc. VLDB Endow.
\textbf{17}(8), 1939--1952 (2024)

\bibitem{wang2021udo}
Wang, J., Trummer, I., Basu, D.: UDO: Universal database optimization
using reinforcement learning. Proc. VLDB Endow. \textbf{14}(13),
3402--3414 (2021)

\end{thebibliography}

\end{document}
